% Owner: classification_e.  This is a restriction on the negative-slope
% obstruction, not an exclusion of all Hopf-surface fibrations.
\subsubsection*{A restriction on the negative-slope fixed divisor}

\begin{proposition}[E.2.8 (Proved): horizontal divisors force non-K\"ahler fibre normalizations]
Let $X$ be a compact complex threefold diffeomorphic to $S^6$, and
let $f:X\to\mathbb P^1$ be a surjective holomorphic map with
connected fibres.  Suppose that a smooth general fibre $F$ has
$b_1(F,\mathbb Q)=1$.  Then every reduced fibre of $f$ is
irreducible, and the monodromy on $H^1(F,\mathbb Q)$ is trivial.

If $H\subset X$ is a horizontal prime divisor, the normalization
of every reduced fibre lies outside Fujiki's class $\mathcal C$.
Equivalently, no such normalization admits a K\"ahler resolution.
In particular no fibre has Moishezon normalization.

Without the hypothesis $b_1(F)=1$, every irreducible component
of a \emph{reducible} fibre of any such connected map from a
complex $S^6$ also has normalization outside $\mathcal C$.

Consequently, in E.2.5 with $D<0$, a horizontal component of
$E_0$ is possible only if every reduced fibre of the central
algebraic reduction $g_0$ is irreducible and has normalization
outside $\mathcal C$.  This necessary condition does not exclude
non-K\"ahler normalizations and does not prove persistence of the
relative pencil in the negative-slope case.
\end{proposition}
\begin{proof}
Choose a nonempty finite set $\Sigma\subset\mathbb P^1$ containing
all critical values; a regular value can be added if necessary.
Write $s=|\Sigma|$, put $A=f^{-1}(\Sigma)_{\mathrm{red}}$,
and let $r$ be the number of irreducible components of $A$.
The compact analytic set $A$ is triangulable.  Each of its
irreducible two-dimensional components has its complex
fundamental class, and their intersections have real dimension
at most two.  Hence
\[
 H^4(A,\mathbb Q)\simeq\mathbb Q^r.
\]
Alexander duality in the oriented six-sphere gives
\[
 b_1(X\setminus A,\mathbb Q)=r.
\]
On the other hand, $f:X\setminus A\to U:=\mathbb P^1\setminus\Sigma$
is a proper smooth fibre bundle.  The base $U$ has the homotopy
type of a graph with $s-1$ loops.  The degree-one Leray--Serre
sequence therefore gives an exact sequence
\[
 0\longrightarrow H^1(U,\mathbb Q)
 \longrightarrow H^1(X\setminus A,\mathbb Q)
 \longrightarrow H^1(F,\mathbb Q)^{\pi_1(U)}
 \longrightarrow0.
\]
In particular, the correct formula retains the invariant
subspace:
\[
 r=s-1+\dim H^1(F,\mathbb Q)^{\pi_1(U)}\le s.
\]
Every nonempty fibre over $\Sigma$ has at least one irreducible
component, so $r\ge s$.  Equality follows throughout.  Thus
every fibre over $\Sigma$ has precisely one reduced irreducible
component and the invariant subspace is all of the
one-dimensional $H^1(F,\mathbb Q)$.  Fibres outside $\Sigma$ are
smooth and connected, hence irreducible.  This proves both
assertions without assuming monodromy triviality at the outset.
It is the rank-one repair of the component-count argument;
the unrestricted formula in \citet[Lemma~4.2]{CDP2020D} is not
being used.

Let $S=f^{-1}(t)_{\mathrm{red}}$.  A horizontal prime $H$ maps
surjectively to the base by properness, so $H\cap S\ne\varnothing$.
Since $S$ is irreducible, $S\subset H$ would imply $S=H$,
contrary to horizontality.  The canonical section of
$\mathcal O_X(H)$ consequently restricts to a nonzero section
on the reduced irreducible space $S$, with a nonempty zero
set.  Its zero divisor is a nonzero effective Cartier divisor
on $S$.  Indeed the section cannot vanish identically on
any local irreducible branch: its pullback would then vanish
on an open subset of the irreducible normalization, hence
everywhere by analytic continuation.  In the reduced local
ring it consequently avoids every minimal prime and is a
non-zero-divisor, including at self-intersections.

Let $\nu:S^\nu\to S$ be the normalization.  If $S^\nu$ were
in $\mathcal C$, it would have a compact K\"ahler resolution
$\mu:\widehat S\to S^\nu$; this precise surface statement is
\citet[Lemma~3.3 and Corollary~3.4]{Fujino2021E}.
The pulled-back line bundle
\[
 L=(\nu\mu)^*(\mathcal O_X(H)|_S)
\]
has $c_1(L)=0$ integrally, since $H^2(X,\mathbb Z)=0$.
Its pulled-back section is nonzero and has zeros: normalization
is finite surjective, and modification is proper surjective.
Thus it defines a nonzero effective divisor $B$ on
$\widehat S$.  For a K\"ahler form $\omega$,
\[
 0=\langle c_1(L)\smile[\omega],[\widehat S]\rangle
   =\int_B\omega>0,
\]
a contradiction.  The argument applies to the reduced support
even when the original fibre is nonreduced.  Moishezon spaces
are in $\mathcal C$, so the stated consequence follows.

For the assertion about a reducible fibre, let $S_i$ be one
of its reduced irreducible components.  The fibre is connected,
so another component $S_j$ meets $S_i$; otherwise $S_i$ would
be an open-and-closed component of its support.  Distinctness
ensures $S_i\not\subset S_j$.  Apply exactly the preceding
canonical-section and positivity argument to
$\mathcal O_X(S_j)|_{S_i}$.  Its first Chern class is zero
because $H^2(X,\mathbb Z)=0$, whereas its effective Cartier
zero divisor is nonempty.  This proves the additional assertion
without any hypothesis on the general fibre's first Betti number.

For the application, choose a very general smooth fibre $F$
of $g_0$ meeting a horizontal component of $E_0$ and avoiding
its vertical components.  As in E.2.6, adjunction gives
$c_1(K_F)=0$ and $F$ is minimal.  The identity of E.2.5 now
restricts to
\[
 K_F^{-A}=\mathcal O_F((-D)E_0|_F),
 \qquad -D>0,
\]
with a nonzero effective divisor on the right.  If a positive
power $K_F^m$ had a nonzero section, multiplying a suitable
power of that section by the $m$th power of the displayed
anticanonical section would give a nonzero holomorphic
function with zeros on $F$.  Compactness forbids this.
Thus $\kappa(F)=-\infty$.  Moreover $F$ is not K\"ahler,
since its nonzero effective divisor $(-D)E_0|_F$ has zero
real Chern class.  The Enriques--Kodaira classification then
places $F$ in class VII and gives $b_1(F)=1$; see
\citet[Chapter~VI, Theorem~1.1 and Table~10, pp.~243--244]
{BHPV2004E}.  The preceding argument applies with $f=g_0$
and any horizontal prime component $H$ of $E_0$.
\end{proof}

The proof supplies a necessary restriction on an unknown
smooth specialization.  It neither classifies the singular
non-K\"ahler normalizations nor replaces the unresolved
non-normal-fibre step in the holomorphic-fibration argument
of \citet[\S7]{CDP2020D}.
